\documentclass{article}
\usepackage[utf8]{inputenc}
\usepackage[left=3cm, right=3cm, top=2cm, bottom=2cm]{geometry}
\usepackage[linesnumbered,boxed]{algorithm2e}
\usepackage[noend]{algpseudocode}
\usepackage{amsmath}
\usepackage{amssymb}
\usepackage{amsthm}
\usepackage{authblk}
\usepackage{comment}
\usepackage{mdframed}
\usepackage{graphicx}

\newtheorem{definition}{Definition}[section]

\begin{document}


\title{Homework 3}
\author{CS 2051, Spring 2022}
\affil{Georgia Tech\\~\\{\bf Due February 12}}
\date{}
\maketitle


\paragraph{Reminder:} You may collaborate with other students in this
class, but (1) you must write up your own solutions in your own words, and (2) 
you must write down everyone you worked with at the top of the page, or ``no 
collaborators'' if you did it all on your own. Additionally, if you used any outside websites or textbooks besides the course text, please cite them here. You may not use question-answer sites like Chegg or MathOverflow. \\

\hline

\subsection*{Question 0}
About how many hours did you spend on this homework in total? Are there any topics that were particularly unclear that we should spend more time on? (not graded)

\subsection*{Question 1}
Let $A, B \subseteq U$ where $U$ is the universe. 
\begin{enumerate}
    \item Prove that $(A \cup B)^C = A^C \cap B^c$.
    % \sam{de morgans law 1 for sets}
    \item Prove that $(A \cap B)^C = A^C \cup B^C$.
    \item Prove that
    \begin{gather*}
            (X \cap Y \cap Z) \cup (X \cup Y \cup Z)^C \\
            \text{and} \\
            (X \cup Y^C) \cap (X \cup Z^C) \cap (Y \cup X^C) \cap (Y \cup Z^C) \cap (Z \cup X^C) \cap (Z \cup Y^C) \\
            \text{are equal.}
    \end{gather*}

    % \sam{ de morgan law 2 for sets}
\end{enumerate}


\subsection*{Question 2}
For each of the following parts, determine if each of the following are true: $A \subseteq B$, $B \subseteq A$, $A = B$. Prove your answers. 
\begin{enumerate}
    \item $A = \{2n : n \in \mathbb{Z} \}, B = \{4n : n \in \mathbb{Z} \}$
    \item $A = \mathbb{Q}, B = \{7n : n \in \mathbb{Q}\}$
    \item $A = \{2n : n \in \mathbb{Z}\}$, $B = \{n = 2k + 1 \text{ for some k in }\mathbb{Z} : n \in \mathbb{R} \}^C$
\end{enumerate}

\subsection*{Question 3}
Out of the following four functions, one is a surjection only, one is an injection only, one is neither, and one is a bijection. State which is which and demonstrate your answers by (1) proving
that it is an injection or naming two distinct elements of the domain that are associated to the
same element of the target, and (2) proving that it is a surjection or naming an element of the
target for which there is no corresponding element in the domain.
\begin{enumerate}
    \item $f: \mathbb{Z}^2 \mapsto \mathbb{Z}$, $f(m, n) = m^2 + n^2$
    \item $f: \mathbb{Z}^2 \mapsto \mathbb{Z}$, $f(m, n) = m$
    \item $f: \mathbb{R}^2 \mapsto \mathbb{C}$, $f(m, n) = m - 2ni$
    \item $f: \mathbb{Z}^2 \mapsto \mathbb{Z}$, $f(m, n) = 2^m3^n$
\end{enumerate}


\subsection*{Question 4}
Determine whether the relation $\mathcal{R}$ on the set of all people is reflexive, symmetric, antisymmetric, and/or transitive. You do not need to justify your answers.\\

$(a,b)\in \mathcal{R}$ if and only if
\begin{enumerate}
\item $a$ is taller than $b$.\\
$\bigcirc$ \texttt{Reflexive} \quad $\bigcirc$ \texttt{Symmetric} \quad $\bigcirc$ \texttt{Antisymmetric} \quad $\bigcirc$ \texttt{Transitive}

\item $a$ and $b$ were born on the same day.\\
$\bigcirc$ \texttt{Reflexive} \quad $\bigcirc$ \texttt{Symmetric} \quad $\bigcirc$ \texttt{Antisymmetric} \quad $\bigcirc$ \texttt{Transitive}

\item $a$ has the same first name as $b$.\\
$\bigcirc$ \texttt{Reflexive} \quad $\bigcirc$ \texttt{Symmetric} \quad $\bigcirc$ \texttt{Antisymmetric} \quad $\bigcirc$ \texttt{Transitive}

\item $a$ and $b$ have a common grandparent.\\
$\bigcirc$ \texttt{Reflexive} \quad $\bigcirc$ \texttt{Symmetric} \quad $\bigcirc$ \texttt{Antisymmetric} \quad $\bigcirc$ \texttt{Transitive}

\item $a$ and $b$ have met.\\ 
$\bigcirc$ \texttt{Reflexive} \quad $\bigcirc$ \texttt{Symmetric} \quad $\bigcirc$ \texttt{Antisymmetric} \quad $\bigcirc$ \texttt{Transitive}

\end{enumerate}

\subsection*{Question 5}
The set of rational numbers can be seen as a set of the equivalence classes of a certain relation on $\mathbb{Z} \times \mathbb{Z}$ defined by $(a, b)R(c, d)$ if and only if $ad = bc$. Show that $R$ is an equivalence relation.


\subsection*{Question 6 \textit{(Hilbert's Grand Hotel)}}
     The Grand Hotel has countable many rooms, numbered 1, 2, 3, \dots, and it is fully occupied.
    \begin{enumerate}
        \item Suppose that $k$ guests arrive at the hotel. Show how you can accommodate all new guests without removing any of the current guests.
        \item Suppose that a countable infinite number of guests arrive. Show how you can accommodate all new guests without removing any of the current guests.
        \item Suppose Hilbert (the owner of the Grand Hotel) expands the property to a second building, also with infinite rooms numbered $1, 2, 3, \dots$. Show how you can spread the current guests to fill every room in both buildings.
    \end{enumerate}

% \subsection*{Solution}
%     \begin{enumerate}
%         \item Shift all existing guests $k$ rooms over, and place the $k$ guests in the rooms $1, \dots k$.
%         \item Move all existing guests in room $n$ to room $2n + 1$. This frees all rooms that can be expressed as $2n$ for the countably infinite number of guests to occupy.
%         \item Take every guest with an odd room number $2k - 1$, where $k \in \mathbb{N}$ and move them into room number $k$ in the second building.
        
%         \vspace{1.5mm}
%         Now take the remaining guests in the original building. Since their room numbers are even, we can represent their room number as $m = 2l$, where $l \in \mathbb{N}$. Move these guests into room number $l$ in the first building.
%     \end{enumerate}

\subsection*{Question 7 (\textit{Cantor's theorem})}



\begin{enumerate}
    \item Prove that there does not exist an onto function between $\mathbb{N}$ and its power set, $\mathcal{P}(\mathbb{N})$.\\

 {\it Hint: proceed by contradiction assuming there is such function $f:\mathbb{N} \rightarrow \mathcal{P}(\mathbb{N})$. Let 
 
 \[
 T=\{n\in \mathbb{N}|~n\notin f(n) \}.
 \]

Since $f$ is onto, there exist $t\in \mathbb{N}$ such that $f(t)=T$. Argue a contradiction by looking at $t\in T$ and $t\notin T$.}
\item Show that the set of all countable sets is uncountable

\end{enumerate}

% There is an easy solution: all sets of the form $\{x\}$ for real values of $x$ are countable, and these is an uncountable subset of the set of countable sets...

% What I would like students to see: all subsets of the natural numbers are countable, but the cardinality of the power set is not countable. Since $\mathcal{P}(\mathbb{N})$ is an uncountable subset of the set of countable sets we are done.



\end{document}